\clearpage
\fancyhead[R]{\small\sffamily\color{Muted}Vulkan / 浴室表面}
\section{浴室的第一层可见表面}
镜头对着洗手台时，一个像素究竟落在镜子、柜门还是后面的墙上？旧的 Vulkan 编辑器从相机出发，给每个像素发一条射线，在场景的三角形中找最近的交点。现在编辑器先把模型投到屏幕，深度测试留下最近的表面；同一遍绘制写出它的底色、法线和深度。现在的两份颜色输出还保存该像素的金属度与粗糙度。接着一次计算把深度还原成世界位置，仍交给原来的直接光和间接光计算。

两份真实编辑器捕获使用同一个 RTXGI 浴室、主相机、$1391\times892$ 的照明画幅，以及第 12 帧的 TAA 采样位置。旧版第 24 号事件的计算参数首项是 0，对应 \code{GI.comp} 查找第一层表面。新版有 85 次绘制，之后的第 459 号事件是 \code{PackSurfaces.comp}，把这两张输出交给后续光照。这里只计这些事件的 GPU 时间，逐次相加后再取中位数。回放一次预热，随后各测 11 次。

\begin{center}\smalltable
\begin{tabular}{lrrr}
\toprule
取得可见表面的方式 & 最短（ms）& 中位数（ms）& 最长（ms）\\
\midrule
相机射线，1 次计算 & 2.281 & 2.435 & 4.152\\
光栅化，85 次绘制＋1 次计算 & 0.375 & 0.391 & 0.752\\
\bottomrule
\end{tabular}
\end{center}

两行完成的着色细节有差别。旧版使用射线命中的三角形几何法线及缩小的贴图，光栅版本可读取材质原图和法线贴图。第二行包含把光栅结果转成表面记录的时间，其中单看第 459 号事件的中位数是 0.147 ms。每笔绘制起止之间的 GPU 时间戳没有覆盖渲染通道开始时的清图、通道结束的存储或事件间的等待。尚未测量相同输入下的 CPU 提交耗时和整帧变化。

下一次测量可以把当前分别提交的表面记录、光照和颜色合成录进同一帧的命令缓冲。已有结果说明可见表面工作本身有多长；另记 CPU 的提交与等待，才知道合并以后实际省掉多少。继续使用这个机位，分别统计静止 32 帧和镜头移动期间的帧时间。

\readings{\href{run:../profiling/primary-surfaces/summary.json}{原始毫秒、文件哈希与事件}\quad
\href{run:../profiling/primary-surfaces/README.md}{RenderDoc 环境与复测命令}\quad
\href{run:Redflag-Source-Guide.pdf}{源码读本 · 石柱的表面与光照}}
